\documentclass[a4paper,12pt]{article}

\usepackage[T1]{fontenc}
\usepackage[utf8]{inputenc}
\usepackage[polish]{babel}
\usepackage{geometry}
\usepackage{graphicx}
\usepackage{booktabs}
\usepackage{longtable}
\usepackage{array}
\usepackage{amsmath}
\usepackage{hyperref}
\usepackage{xcolor}

\geometry{margin=2.5cm}

\title{Radiowy Translator Kodu Morse'a}
\author{
AA\\nr indeksu 111
\and
BB\\nr indeksu 222
\and
CC (Kierownik Projektu)\\nr indeksu 333
}
\date{\today}

\begin{document}

\maketitle

\tableofcontents
\newpage

\section{Cel projektu}

Celem projektu jest realizacja systemu transmisji informacji wykorzystującego kod Morse'a oraz komunikację radiową w paśmie 433.92 MHz.

System składa się z dwóch urządzeń:

\begin{itemize}
\item nadajnika,
\item odbiornika.
\end{itemize}

Nadajnik umożliwia wprowadzanie informacji za pomocą przycisku. Informacja jest interpretowana jako kod Morse'a, następnie kodowana i przesyłana drogą radiową.

Odbiornik odbiera komunikat, dekoduje go i prezentuje użytkownikowi w postaci tekstowej oraz dźwiękowej.

\section{Lista funkcjonalności}

\begin{enumerate}
\item Wprowadzanie znaków kodu Morse'a za pomocą przycisku.
\item Rozróżnianie kropek i kresek na podstawie czasu przytrzymania.
\item Dekodowanie kodu Morse'a do znaków alfabetu łacińskiego.
\item Sygnalizacja świetlna poprawnego przetłumaczenia ciągu znaków.
\item Regulacja szybkości interpretacji kodu Morse'a.
\item Transmisja radiowa 433 MHz pomiędzy urządzeniami.
\item Odbiór komunikatów radiowych.
\item Wyświetlanie odebranego komunikatu na wyświetlaczu LCD 16x2.
\item Sygnalizacja akustyczna odebrania komunikatu.
\item Sygnalizacja świetlna poprawnego odbioru komunikatu.
\item Zliczanie odebranych komunikatów.
\end{enumerate}

\section{Podział obowiązków}

\begin{table}[h]
\centering
\begin{tabular}{|p{4cm}|p{2cm}|p{7cm}|}
\hline
Osoba & Udział & Zakres prac \\
\hline
AA & 40\% &
Projekt sprzętu, montaż układów, testy komunikacji radiowej \\
\hline
BB & 30\% &
Obsługa LCD, dokumentacja, analiza FMEA \\
\hline
CC & 30\% &
Implementacja algorytmów, dekodowanie Morse'a, integracja systemu \\
\hline
\end{tabular}
\end{table}

\section{Instrukcja użytkownika}

\begin{enumerate}
\item Włączyć nadajnik i odbiornik.
\item Odczekać zakończenie inicjalizacji urządzeń.
\item Wprowadzać kod Morse'a:
\begin{itemize}
\item naciśnięcie krótsze niż 0.3 s oznacza kropkę,
\item naciśnięcie dłuższe niż 0.3 s oznacza kreskę.
\end{itemize}
\item Po zakończeniu litery odczekać około 0.8 s.
\item Znak zostanie przesłany drogą radiową.
\item Odebrany znak zostanie wyświetlony na LCD.
\item Poprawny odbiór zostanie zasygnalizowany diodą LED i buzzerem.
\item Przyciski SPEED+ i SPEED− umożliwiają zmianę szybkości interpretacji kodu Morse'a.
\end{enumerate}

\section{Opis sprzętu}

\subsection{Nadajnik}

Mikrokontroler:
\begin{itemize}
\item ATmega328PB (Arduino Nano v3)
\end{itemize}

Elementy:
\begin{itemize}
\item przycisk Morse'a,
\item przycisk zwiększania szybkości,
\item przycisk zmniejszania szybkości,
\item dioda LED,
\item buzzer,
\item moduł RF FS1000A.
\end{itemize}

\subsection{Odbiornik}

Mikrokontroler:
\begin{itemize}
\item ATmega328PB (Arduino Nano v3)
\end{itemize}

Elementy:
\begin{itemize}
\item odbiornik RF XY-MK,
\item wyświetlacz LCD1602 z interfejsem I2C,
\item dioda LED,
\item buzzer.
\end{itemize}

\section{Mapowanie funkcjonalności na sprzęt}

\begin{table}[h]
\centering
\begin{tabular}{|l|l|}
\hline
Funkcja & Element sprzętowy \\
\hline
Wprowadzanie Morse'a & Przycisk D4 \\
Regulacja szybkości & SPEED+, SPEED− \\
Sygnalizacja lokalna & LED + buzzer \\
Transmisja danych & FS1000A \\
Odbiór danych & XY-MK \\
Wyświetlanie danych & LCD1602 \\
\hline
\end{tabular}
\end{table}

\section{Opis działania algorytmu}

System oczekuje na naciśnięcie przycisku.

W chwili naciśnięcia zapamiętywany jest czas:

\[
t_0
\]

W chwili zwolnienia przycisku zapamiętywany jest czas:

\[
t_1
\]

Obliczana jest długość impulsu:

\[
\\Delta t = t_1 - t_0
\]

Jeżeli:

\[
\\Delta t < T_m
\]

impuls interpretowany jest jako kropka.

Jeżeli:

\[
\\Delta t \ge T_m
\]

impuls interpretowany jest jako kreska.

Symbol zostaje zapisany do bufora kodu Morse'a.

Jeżeli czas bezczynności przekroczy 0.8 s, następuje dekodowanie ciągu na znak ASCII.

Znak zostaje przesłany drogą radiową.

Odbiornik odbiera komunikat i prezentuje go użytkownikowi za pomocą wyświetlacza LCD, diody LED oraz buzzera.

\section{Parametry czasowe}

\begin{table}[h]
\centering
\begin{tabular}{|l|l|}
\hline
Parametr & Wartość \\
\hline
Próg kropka/kreska & 0.3 s \\
Przerwa kończąca literę & 0.8 s \\
Napięcie zasilania & 5 V \\
Częstotliwość MCU & 16 MHz \\
Częstotliwość RF & 433.92 MHz \\
\hline
\end{tabular}
\end{table}

\section{Opis peryferiów i rejestrów}

\subsection{GPIO}

Do obsługi wejść i wyjść cyfrowych wykorzystano rejestry:

\begin{itemize}
\item DDRD -- konfiguracja kierunku portu,
\item PORTD -- ustawienie stanu logicznego oraz aktywacja rezystorów podciągających,
\item PIND -- odczyt aktualnego stanu wejścia.
\end{itemize}

Przykład konfiguracji:

\begin{verbatim}
DDRD |= (1<<DDD6);
\end{verbatim}

Powoduje ustawienie pinu PD6 jako wyjścia.

\subsection{Timer0}

Pomiar czasu realizowany jest przy pomocy sprzętowego licznika Timer0.

Parametry:

\begin{itemize}
\item częstotliwość zegara: 16 MHz,
\item preskaler: 64,
\item rozdzielczość funkcji millis(): 1 ms.
\end{itemize}

\subsection{I2C (TWI)}

Do komunikacji z wyświetlaczem LCD wykorzystano interfejs TWI.

Wykorzystywane rejestry:

\begin{itemize}
\item TWBR -- ustalenie prędkości transmisji,
\item TWCR -- sterowanie transmisją,
\item TWDR -- rejestr danych,
\item TWSR -- rejestr statusu.
\end{itemize}

Prędkość transmisji:

[
100,kHz
]

Adres urządzenia:

[
0x27
]

\section{Komunikacja radiowa}

Komunikacja realizowana jest przy użyciu transmisji OOK (On-Off Keying) w paśmie 433.92 MHz.

Moduł FS1000A generuje sygnał radiowy odpowiadający transmitowanym bitom.

Stan logiczny 1 reprezentowany jest obecnością nośnej.

Stan logiczny 0 reprezentowany jest brakiem nośnej.

Odbiornik XY-MK demoduluje sygnał i odtwarza przebieg cyfrowy.

\section{FMEA}

\begin{table}[h]
\centering
\begin{tabular}{|p{4cm}|c|c|c|c|}
\hline
Ryzyko & P & S & D & P$\times$S$\times$D \\
\hline
Zablokowany przycisk & 0.2 & 6 & 8 & 9.6 \\
\hline
Zakłócenia RF & 0.4 & 8 & 5 & 16 \\
\hline
Nieznany kod Morse'a & 0.3 & 3 & 10 & 9 \\
\hline
Przepełnienie LCD & 0.5 & 2 & 10 & 10 \\
\hline
\end{tabular}
\end{table}

Najbardziej krytycznym zagrożeniem są zakłócenia transmisji radiowej.

\section{Dokumentacja kodu}

Każda funkcja programu opisywana jest zgodnie ze standardem Doxygen.

Przykład:

\begin{verbatim}
/*!

* @brief Dekodowanie ciągu Morse'a.
*
* @param code Ciąg znaków '.' i '-'.
*
* @return Znak ASCII lub '?'.
  */
  char DecodeMorse(const char* code);
  \end{verbatim}

\section{Bibliografia}

\begin{enumerate}
\item ATmega328P Datasheet.
\item ATmega328PB Datasheet.
\item HD44780 LCD Controller Datasheet.
\item PCF8574 I/O Expander Datasheet.
\item I2C Bus Specification.
\item FS1000A RF Transmitter Documentation.
\item XY-MK RF Receiver Documentation.
\end{enumerate}

\section{Załączniki}

\begin{itemize}
\item Schemat nadajnika.
\item Schemat odbiornika.
\item Kod źródłowy.
\item Plik wykonywalny HEX.
\item Noty katalogowe wykorzystanych elementów.
\end{itemize}

\end{document}
